\documentclass[9pt,a4paper]{extarticle}

% ======================
% Codificação, idioma e layout
% ======================
\usepackage[T1]{fontenc}
\usepackage[utf8]{inputenc}
\usepackage[brazil]{babel}
\usepackage{geometry}
\geometry{margin=2cm}
\usepackage{parskip}
\usepackage{setspace}
\usepackage{microtype}
\usepackage{enumitem}
\usepackage{graphicx}
\usepackage{algpseudocode}
\usepackage{color}

% ======================
% Matemática e símbolos
% ======================
\usepackage{amsmath, amssymb, amsthm, bm}

% ======================
% Figuras e tabelas
% ======================
\usepackage{graphicx}
\usepackage{booktabs}
\usepackage{array}
\usepackage{tabularx}
\usepackage{float}
\usepackage{caption}
\usepackage{subcaption}

% ======================
% Hiperlinks
% ======================
\usepackage{hyperref}
\hypersetup{
  colorlinks=true,
  linkcolor=blue,
  urlcolor=blue
}

% =====================================================================
% Solução — Análise de Busca-Linear e Tem-Duplicata
% Paleta: azulescuro = RGB (10, 40, 90); laranja = RGB (220, 110, 30)
% =====================================================================
\definecolor{azulescuro}{RGB}{10,40,90}
\definecolor{laranja}{RGB}{220,110,30}

% ======================
% Início do documento
% ======================
\begin{document}

\begin{center}
  \Large \textbf{Exercícios Resolvidos - Lista 2 - Análise de Algoritmos e Caracterização de Tempos de Execução} \\[4pt]
  \normalsize Disciplina: PCC104 - Projeto de Análise de Algoritmos
\end{center}


% --------------------------------------------------------------------

\begin{itemize}
    \item \textbf{Exercício 2}


\begin{enumerate}[label=(\alph*)]
    \item \textbf{Tamanho da entrada.} O tamanho da entrada é $n$, o número de elementos do arranjo $A$. O valor $x$ e cada $A[i]$ ocupam $\Theta(1)$ espaço no modelo RAM usual (palavras de tamanho fixo), de modo que o tamanho dominante é o comprimento do arranjo.

    \item \textbf{Tabela de custos.} Seja $t$ o número de vezes que o teste do \textbf{while} na linha~2 é avaliado. Como o teste é avaliado uma vez extra (aquela que falha e encerra o laço), a linha~3 (incremento) executa exatamente $t-1$ vezes.

    \begin{center}
    \renewcommand{\arraystretch}{1.15}
    \begin{tabular}{@{}clcc@{}}
        \toprule
        \textbf{Linha} & \textbf{Instrução} & \textbf{Custo} & \textbf{Vezes} \\
        \midrule
        1 & $i \gets 1$                              & $c_1$ & $1$    \\
        2 & \textbf{while} $i \leq n$ \textbf{and} $A[i]\neq x$ & $c_2$ & $t$    \\
        3 & \quad $i \gets i+1$                      & $c_3$ & $t-1$  \\
        4 & \textbf{if} $i \leq n$                    & $c_4$ & $1$    \\
        5 & \quad \textbf{return} $i$                 & $c_5$ & $0$ ou $1$ \\
        7 & \quad \textbf{return} $-1$                & $c_7$ & $0$ ou $1$ \\
        \bottomrule
    \end{tabular}
    \end{center}

    Exatamente uma das linhas~5 ou~7 executa; denotemos por $c_{\text{ret}}$ esse custo terminal (constante, independente de $n$).

    \item \textbf{Expressão de $T(n)$.} Somando custo $\times$ frequência:
    \[
        T(n) \;=\; c_1 + c_2\,t + c_3\,(t-1) + c_4 + c_{\text{ret}}
              \;=\; (c_2 + c_3)\,t \;+\; \underbrace{(c_1 + c_4 + c_{\text{ret}} - c_3)}_{\text{constante}}.
    \]

    \item \textbf{\textcolor{laranja}{Melhor caso.}} Ocorre quando o teste falha já na primeira iteração porque $A[1] = x$. Nesse caso, o \textbf{while} avalia o teste uma única vez (e sai), logo $t = 1$.

    \emph{Instância:} qualquer arranjo com $A[1] = x$, por exemplo $A = \langle x, *, *, \ldots, *\rangle$.

    Substituindo $t = 1$:
    \[
        T_{\text{melhor}}(n) = (c_2 + c_3)\cdot 1 + (c_1 + c_4 + c_{\text{ret}} - c_3) = c_1 + c_2 + c_4 + c_{\text{ret}},
    \]
    que é uma constante. Logo $T_{\text{melhor}}(n) = \Theta(1)$, com constantes $c_1 = c_2 = c_1 + c_2 + c_4 + c_{\text{ret}}$ e $n_0 = 1$ (na definição $c_1 \leq T_{\text{melhor}}(n) \leq c_2$ para todo $n \geq n_0$).

    \item \textbf{\textcolor{laranja}{Pior caso.}} Ocorre quando o laço percorre todas as posições sem nunca encontrar $x$: o teste avalia para $i = 1, 2, \ldots, n$ (verdadeiro) e mais uma vez para $i = n+1$ (falso), totalizando $t = n+1$.

    \emph{Instâncias:} todos os arranjos $A$ tais que $A[i] \neq x$ para todo $i \in \{1, \ldots, n\}$, isto é, $x \notin A$.

    Substituindo $t = n+1$:
    \[
        T_{\text{pior}}(n) = (c_2 + c_3)(n+1) + (c_1 + c_4 + c_{\text{ret}} - c_3) = (c_2 + c_3)\,n + \kappa,
    \]
    com $\kappa$ constante. Tomando $a = c_2 + c_3 > 0$, vale $a n \leq T_{\text{pior}}(n) \leq 2 a n$ para $n \geq n_0$ suficientemente grande (basta $n_0 \geq |\kappa|/a$), portanto $T_{\text{pior}}(n) = \Theta(n)$.

    \item \textbf{\textcolor{laranja}{Caso médio.}} Seja $X$ a variável aleatória ``posição da única ocorrência de $x$ em $A$'', uniforme em $\{1, 2, \ldots, n\}$. Se $X = k$, o teste do \textbf{while} é avaliado $k$ vezes: $k-1$ vezes com resultado verdadeiro (linhas $A[i] \neq x$ para $i < k$) e uma vez com resultado falso em $i = k$. Logo $t \mid (X = k) = k$.

    \[
        \mathbb{E}[t] \;=\; \sum_{k=1}^{n} k \cdot \Pr[X = k] \;=\; \frac{1}{n}\sum_{k=1}^{n} k \;=\; \frac{1}{n} \cdot \frac{n(n+1)}{2} \;=\; \frac{n+1}{2}.
    \]

    Por linearidade da esperança aplicada a $T(n) = (c_2 + c_3)\, t + \kappa'$:
    \[
        T_{\text{med}}(n) \;=\; \mathbb{E}[T(n)] \;=\; (c_2 + c_3)\,\frac{n+1}{2} + \kappa' \;=\; \tfrac{c_2 + c_3}{2}\, n + \kappa'',
    \]
    com $\kappa''$ constante. Concluímos que $T_{\text{med}}(n) = \Theta(n)$.

    \emph{Observação:} o caso médio cresce na mesma classe do pior caso, mas com constante multiplicativa exatamente a metade.
\end{enumerate}

% ---------------------------------------------------------------------
\item \textbf{Exercício 3}

\begin{enumerate}[label=(\alph*)]
    \item \textbf{Tamanho da entrada.} O tamanho da entrada é $n$, o número de elementos do arranjo $A$.

    \item \textbf{Pior caso: contagem do laço interno.} Para cada valor de $i \in \{1, \ldots, n-1\}$, o laço interno itera $j$ de $i+1$ até $n$, ou seja, $n - i$ iterações; em cada uma delas o teste da linha~2 é avaliado uma vez (a avaliação extra de saída custa apenas o teste falso da condição do \textbf{for}, que aqui contabilizamos junto). No pior caso, o algoritmo nunca retorna na linha~4, percorrendo todos os pares.

    O número total de avaliações da comparação interna é
    \[
        \sum_{i=1}^{n-1} (n - i) \;=\; \frac{n(n-1)}{2}.
    \]

    \emph{Instância:} qualquer arranjo com todos os elementos distintos, por exemplo $A = \langle 1, 2, 3, \ldots, n\rangle$.

    \item \textbf{Tabela de custos --- pior caso.} Seja $S = \dfrac{n(n-1)}{2}$ o número de pares $(i,j)$ com $i < j$.

    \begin{center}
    \renewcommand{\arraystretch}{1.15}
    \begin{tabular}{@{}clcc@{}}
        \toprule
        \textbf{Linha} & \textbf{Instrução} & \textbf{Custo} & \textbf{Vezes} \\
        \midrule
        1 & \textbf{for} $i \gets 1$ \textbf{to} $n-1$               & $c_1$ & $n$ \\
        2 & \quad \textbf{for} $j \gets i+1$ \textbf{to} $n$          & $c_2$ & $\sum_{i=1}^{n-1}(n-i+1) = S + (n-1)$ \\
        3 & \quad\quad \textbf{if} $A[i] = A[j]$                      & $c_3$ & $S$ \\
        4 & \quad\quad\quad \textbf{return verdadeiro}                & $c_4$ & $0$ \\
        5 & \textbf{return falso}                                      & $c_5$ & $1$ \\
        \bottomrule
    \end{tabular}
    \end{center}

    Observação: para o laço interno na linha~2, cada bloco de $i$ produz $(n - i)$ testes verdadeiros mais $1$ teste falso (que encerra o laço), totalizando $n - i + 1$ avaliações; somando em $i$ obtém-se $S + (n-1)$. A linha~3 (comparação) é executada apenas dentro do laço, $S$ vezes.

    Portanto,
    \[
        T_{\text{pior}}(n) \;=\; c_1\,n + c_2\,\bigl(S + (n-1)\bigr) + c_3\,S + c_5,
    \]
    e substituindo $S = \frac{n(n-1)}{2}$:
    \[
        T_{\text{pior}}(n) \;=\; \frac{c_2 + c_3}{2}\, n^2 \;+\; \left(c_1 + \frac{c_2 - c_3}{2}\right) n \;+\; (c_5 - c_2).
    \]

    \item \textbf{$T_{\text{pior}}(n) = \Theta(n^2)$.} Escrevemos $T_{\text{pior}}(n) = a\,n^2 + b\,n + d$ com
    \[
        a = \frac{c_2 + c_3}{2} > 0, \qquad b = c_1 + \frac{c_2 - c_3}{2}, \qquad d = c_5 - c_2.
    \]
    Tome as constantes $\alpha_1 = a/2$, $\alpha_2 = 2a$ e $n_0$ grande o suficiente para que
    \[
        \alpha_1 n^2 \;\leq\; a\,n^2 + b\,n + d \;\leq\; \alpha_2 n^2 \qquad \text{para todo } n \geq n_0.
    \]
    Para a cota superior, basta $n_0$ tal que $b\,n + d \leq a\,n^2$ a partir dele (válido pois o termo quadrático domina). Para a cota inferior, basta $n_0$ tal que $b\,n + d \geq -\tfrac{a}{2}\,n^2$ a partir dele. Ambos os $n_0$ são finitos, logo $T_{\text{pior}}(n) = \Theta(n^2)$.

    \item \textbf{\textcolor{laranja}{Melhor caso.}} Ocorre quando $A[1] = A[2]$: na primeira iteração do laço externo, com $i = 1$ e $j = 2$, a linha~3 executa uma única vez, o teste é verdadeiro e o algoritmo retorna na linha~4.

    Número de execuções da linha~3: exatamente $1$. Cada uma das demais linhas executa $O(1)$ vezes. Portanto
    \[
        T_{\text{melhor}}(n) = \Theta(1).
    \]

    \item \textbf{Comparação com \textsc{Insertion-Sort}.} Tanto \textsc{Tem-Duplicata} quanto \textsc{Insertion-Sort} apresentam $T_{\text{pior}}(n) = \Theta(n^2)$. A semelhança é estrutural: ambos os algoritmos varrem todos (ou quase todos) os pares de índices $(i, j)$ com $i < j$, o que naturalmente leva à soma $\sum_{i=1}^{n-1}(n-i) = \frac{n(n-1)}{2}$. A diferença está no melhor caso: \textsc{Insertion-Sort} sobre um arranjo já ordenado executa apenas a comparação inicial de cada iteração externa, atingindo $\Theta(n)$; \textsc{Tem-Duplicata}, por sua vez, sai imediatamente quando o primeiro par $(1, 2)$ já é igual, atingindo $\Theta(1)$.
\end{enumerate}

\item \textbf{Exercício 10}

% =====================================================================
% Solução --- MÁXIMO por divisão e conquista
% Paleta: azulescuro = RGB (10, 40, 90); laranja = RGB (220, 110, 30)
% =====================================================================

\begin{enumerate}[label=(\alph*)]
    \item \textbf{Passos do paradigma D\&C.}
    \begin{itemize}
        \item \textbf{Dividir:} calcular $q = \lfloor (p + r)/2 \rfloor$ (linha~4) particiona $A[p\,..\,r]$ em dois subarranjos contíguos $A[p\,..\,q]$ e $A[q+1\,..\,r]$, de tamanhos $\lceil n/2 \rceil$ e $\lfloor n/2 \rfloor$. Custa $\Theta(1)$.
        \item \textbf{Conquistar:} as chamadas recursivas das linhas~5 e~6 resolvem os dois subproblemas independentemente, retornando os máximos $m_1$ e $m_2$. O caso-base é $p = r$ (linhas~1--3), em que o máximo é o próprio elemento $A[p]$.
        \item \textbf{Combinar:} a comparação $m_1 \geq m_2$ e a devolução do maior valor (linhas~7--11) combinam as duas soluções parciais em $\Theta(1)$.
    \end{itemize}

    \item \textbf{Relação de recorrência.} Cada chamada não trivial dispara $2$ chamadas recursivas, cada uma sobre um subarranjo de tamanho aproximadamente $n/2$; o trabalho fora das chamadas recursivas (divisão $+$ combinação) é $\Theta(1)$. O caso-base $n = 1$ executa apenas o teste da linha~1 e o retorno da linha~2, custando $\Theta(1)$. Logo,
    \[
        T(n) \;=\;
        \begin{cases}
            \Theta(1), & \text{se } n = 1, \\[2pt]
            2\,T(n/2) + \Theta(1), & \text{se } n > 1.
        \end{cases}
    \]
    Para simplificar a álgebra, escrevemos $T(n) = 2\,T(n/2) + c$ com $c > 0$ constante e $T(1) = c$.

    \item \textbf{\textcolor{laranja}{Árvore de recursão, $n = 2^k$.}} A cada nível, cada subproblema gera $2$ filhos com metade do tamanho:

    \begin{center}
    \renewcommand{\arraystretch}{1.2}
    \begin{tabular}{@{}cccc@{}}
        \toprule
        \textbf{Nível $\ell$} & \textbf{Nós} & \textbf{Tamanho de cada subproblema} & \textbf{Custo do nível} \\
        \midrule
        $0$        & $1$            & $n$       & $c$            \\
        $1$        & $2$            & $n/2$     & $2c$           \\
        $2$        & $4$            & $n/4$     & $4c$           \\
        $\vdots$   & $\vdots$       & $\vdots$  & $\vdots$       \\
        $\ell$     & $2^\ell$       & $n/2^\ell$ & $2^\ell\, c$   \\
        $\vdots$   & $\vdots$       & $\vdots$  & $\vdots$       \\
        $k = \log_2 n$ & $2^k = n$  & $1$       & $n \cdot c$    \\
        \bottomrule
    \end{tabular}
    \end{center}

    A árvore tem $k + 1 = \log_2 n + 1$ níveis. O custo total é a soma sobre todos os níveis:
    \[
        T(n) \;=\; \sum_{\ell=0}^{k} 2^\ell\, c \;=\; c\,(2^{k+1} - 1) \;=\; c\,(2n - 1).
    \]
    Note que, ao contrário do que acontece no \textsc{Merge-Sort}, aqui o custo por nível \emph{não} é constante: ele dobra a cada nível porque $f(n) = \Theta(1)$ e o número de subproblemas cresce geometricamente. Em consequência, o custo é dominado pelas folhas.

    \item \textbf{$T(n) = \Theta(n)$.} De $T(n) = c(2n - 1) = 2c\,n - c$, tomemos $\alpha_1 = c$ e $\alpha_2 = 2c$. Para todo $n \geq 1$ vale
    \[
        \alpha_1\, n \;=\; c\,n \;\leq\; 2c\,n - c \;\leq\; 2c\,n \;=\; \alpha_2\, n,
    \]
    pois $c\,n \leq 2c\,n - c \iff c \leq c\,n$, válido para $n \geq 1$. Portanto $T(n) = \Theta(n)$.

    \item \textbf{\textcolor{laranja}{Teorema Mestre.}} Comparando $T(n) = 2\,T(n/2) + \Theta(1)$ com a forma canônica $T(n) = a\,T(n/b) + f(n)$:
    \[
        a = 2, \qquad b = 2, \qquad f(n) = \Theta(1).
    \]
    O expoente crítico é $\log_b a = \log_2 2 = 1$, logo $n^{\log_b a} = n$. Comparando $f(n)$ com $n^{\log_b a}$:
    \[
        f(n) = \Theta(1) = O(n^{1 - \varepsilon}) \quad \text{para qualquer } 0 < \varepsilon \leq 1.
    \]
    Estamos no \textbf{Caso~1} do Teorema Mestre, que conclui
    \[
        T(n) = \Theta(n^{\log_b a}) = \Theta(n),
    \]
    confirmando o resultado obtido pela árvore de recursão.

    \item \textbf{Comparação com a versão iterativa.} A varredura linear --- inicializar $\max \gets A[1]$ e percorrer $A[2\,..\,n]$ atualizando $\max$ quando $A[i] > \max$ --- também é $\Theta(n)$. As duas abordagens estão na mesma classe assintótica.

    Há, porém, uma diferença nas constantes e no consumo de memória. A versão recursiva executa $2n - 1$ chamadas no total ($n$ folhas $+$ $n-1$ nós internos) e usa pilha de profundidade $\Theta(\log n)$; a iterativa faz exatamente $n - 1$ comparações com $\Theta(1)$ de espaço auxiliar.

    A lição é que divisão e conquista \emph{não traz ganho assintótico} quando o passo de combinação é $\Theta(1)$ e o problema original já admite solução linear: nesse cenário, $T(n) = 2T(n/2) + \Theta(1)$ cai no Caso~1 do Mestre e o custo é dominado pelas folhas, recuperando o mesmo $\Theta(n)$ da varredura. O paradigma só compensa quando combinar é mais barato que resolver do zero --- como em \textsc{Merge-Sort}, onde combinar custa $\Theta(n)$ mas resolver iterativamente custaria $\Theta(n^2)$ (\textsc{Insertion-Sort}).
\end{enumerate}

\item \textbf{Exercício 9}

\begin{enumerate}[label=(\alph*)]
    \item \textbf{Tamanho dos subproblemas.}
    
    Seja $n = r - p + 1$ o tamanho do subarranjo corrente e $q = \lfloor (p+r)/2 \rfloor$. As duas chamadas recursivas operam sobre os subarranjos $A[p\,..\,q-1]$ e $A[q+1\,..\,r]$, cujos tamanhos são, respectivamente,
    \[
        n_{\text{esq}} = (q-1) - p + 1 = q - p, 
        \qquad
        n_{\text{dir}} = r - (q+1) + 1 = r - q.
    \]
    Substituindo $q = \lfloor (p+r)/2 \rfloor$:
    \[
        n_{\text{esq}} = \left\lfloor \frac{p+r}{2} \right\rfloor - p 
                       = \left\lfloor \frac{r-p}{2} \right\rfloor
                       = \left\lfloor \frac{n-1}{2} \right\rfloor 
                       \leq \left\lfloor \frac{n}{2} \right\rfloor,
    \]
    \[
        n_{\text{dir}} = r - \left\lfloor \frac{p+r}{2} \right\rfloor 
                       = \left\lceil \frac{r-p}{2} \right\rceil 
                       = \left\lceil \frac{n-1}{2} \right\rceil 
                       \leq \left\lfloor \frac{n}{2} \right\rfloor.
    \]
    A última desigualdade vale porque $\lceil (n-1)/2 \rceil = \lfloor n/2 \rfloor$ para todo $n \geq 1$. Assim, qualquer que seja o ramo escolhido, o subproblema tem tamanho \emph{no máximo} $\lfloor n/2 \rfloor$.
    
    \medskip
    
    \item \textbf{Recorrência do pior caso.}
    
    O pior caso ocorre quando $x \notin A$ ou quando $x$ se encontra em uma posição que exige descer até um subarranjo unitário. Em cada chamada, o trabalho não recursivo (cálculo de $q$, comparações nas linhas 1, 4 e 6 e retorno) é $\Theta(1)$. Há \emph{uma} chamada recursiva sobre um subproblema de tamanho no máximo $\lfloor n/2 \rfloor$. O caso base ocorre quando $p > r$, isto é, $n = 0$, com custo constante. Portanto,
    \[
        T(n) = 
        \begin{cases}
            \Theta(1), & \text{se } n \leq 1, \\[4pt]
            T(\lfloor n/2 \rfloor) + \Theta(1), & \text{se } n \geq 2.
        \end{cases}
    \]
    
    \medskip
    
    \item \textbf{Resolução por iteração para $n = 2^k$.}
    
    Para $n = 2^k$, os pisos são exatos: $\lfloor n/2 \rfloor = n/2 = 2^{k-1}$. Seja $c$ uma constante positiva que limita o custo não recursivo. Então:
    \begin{align*}
        T(n) &= T(n/2) + c \\
             &= \bigl[T(n/4) + c\bigr] + c = T(n/4) + 2c \\
             &= \bigl[T(n/8) + c\bigr] + 2c = T(n/8) + 3c \\
             &\;\;\vdots \\
             &= T(n/2^i) + i\,c.
    \end{align*}
    A iteração termina quando $n/2^i = 1$, ou seja, $i = \log_2 n = k$. Logo,
    \[
        T(n) = T(1) + k\,c = \Theta(1) + c \log_2 n.
    \]
    Tomando $T(1) = d$ constante, obtemos a forma fechada
    \[
        \boxed{\,T(n) = c \log_2 n + d.\,}
    \]
    
    \medskip
    
    \item \textbf{Classe assintótica: $T(n) = \Theta(\log n)$.}
    
    Da expressão fechada $T(n) = c \log_2 n + d$, exibimos constantes $c_1, c_2 > 0$ e $n_0$ tais que
    \[
        c_1 \log_2 n \leq T(n) \leq c_2 \log_2 n, \qquad \forall\, n \geq n_0.
    \]
    Basta tomar, por exemplo, $n_0 = 2$, $c_1 = c$ (pois $d \geq 0$) e $c_2 = c + d$ (pois $\log_2 n \geq 1$ para $n \geq 2$, donde $d \leq d \log_2 n$). Portanto $T(n) = \Theta(\log_2 n) = \Theta(\log n)$, já que logaritmos em bases distintas diferem apenas por fator constante.
    
    \medskip
    
    \item \textbf{Melhor caso.}
    
    O melhor caso ocorre quando $A[q] = x$ na \emph{primeira} chamada, isto é, quando o elemento procurado encontra-se exatamente na posição mediana inicial. Nesse cenário não há recursão e o algoritmo executa apenas trabalho constante:
    \[
        T_{\text{melhor}}(n) = \Theta(1).
    \]
    
    \medskip
    
    \item \textbf{Aplicação do Teorema Mestre.}
    
    A recorrência do pior caso tem a forma $T(n) = a\,T(n/b) + f(n)$, com
    \[
        a = 1, \qquad b = 2, \qquad f(n) = \Theta(1).
    \]
    Calculamos o expoente crítico:
    \[
        \log_b a = \log_2 1 = 0, 
        \qquad 
        n^{\log_b a} = n^0 = 1.
    \]
    Comparando $f(n) = \Theta(1)$ com $n^{\log_b a} = \Theta(1)$, vemos que ambos são polinomialmente equivalentes, ou seja, $f(n) = \Theta(n^{\log_b a})$. Aplica-se, portanto, o \textbf{Caso 2} do Teorema Mestre, que fornece
    \[
        T(n) = \Theta\!\left(n^{\log_b a} \log n\right) = \Theta(\log n),
    \]
    confirmando o resultado obtido pelo método da iteração.
\end{enumerate}

\end{itemize}

\end{document}